\clearpage
\fancyhead[L]{\small MERCURY TABLES.}
\SourceLeaf{564}
\begin{center}\small Göttingische gelehrte Anzeigen. No. 107. Pages 1065..1069. 5 July 1813.\\[3pt]\rule{.72\linewidth}{.3pt}\end{center}

Mr. \textsc{von Lindenau}, a corresponding member of the Royal Society of Sciences, has already made a valuable contribution to planetary theory with his tables of Venus and Mars. In a letter to Professor \textsc{Gauss} dated 20 June of this year, he gave a preliminary account of a similar and no less important third work, on \emph{Mercury}, which he recently completed. The worthy efforts of \textsc{Lalande}, \textsc{Oriani}, and \textsc{Triesnecker} had already brought the theory of Mercury's orbit to a high degree of perfection. Even so, practical and theoretical aids were available, some not yet used and some newly available. It was therefore not unlikely that a fresh investigation, carefully taking account of everything useful that astronomy now offers for this purpose, would yield further changes and improvements to the existing determinations. This prospect, together with the hope that a precise comparison of the perturbations of Mercury's elliptical orbit given by theory with those inferred from observations would yield a new and reliable determination of the mass of Venus, was what chiefly prompted Mr. \textsc{von Lindenau} to take up the subject afresh. We give only a few of the final results here. The details of these investigations will be presented to astronomers in a work due to appear in print in a few weeks under the title \emph{Investigatio nova orbitae a Mercurio circa Solem descriptae: accedunt tabulae Planetae ex elementis recens correctis et ex theoria gravitatis ill. de la Place constructae, auctore} \textsc{Bernhardo de Lindenau}.

In keeping with the purpose indicated above, the investigation falls into four sections. The first two discuss the transits of Mercury observed since 1631 and use them to determine the orbit. The latter sections do the same with more recent geocentric observations. Two procedures were applied to the transits of Mercury. First, the node, the motion of the node, and the inclination could be determined from them. Second, a distinctive feature of Mercury's orbit is that all transits at the same node take place at nearly the same points of the orbit. This made it possible to reduce the heliocentric longitudes in the ecliptic given by the transits to mean longitudes in the orbit without any elements other than the mean annual motion and the annual change of the aphelion being able to affect the difference between two such longitudes and make it substantially inaccurate. Applying this procedure does, however, require prior approximate knowledge of Mercury's orbit, which was indeed available in this case. Only twelve transits could be used for the first determination, which chiefly required accurate knowledge of the geocentric latitude. From the seventeen transits discussed and calculated by the author in all, fifteen combinations were formed that seemed to him best suited to the numerical development and application of the second procedure. The annual changes in the node and aphelion obtained in this way were compared with those given by theory, whose values depend chiefly on the mass of Venus. This supplied the first two equations for determining that mass.

The heliocentric longitudes of Mercury derived from the seventeen transits just mentioned, together with a hundred geocentric longitudes observed by \textsc{Maskelyne} and \textsc{Piazzi}, were used to determine the strictly elliptical elements of Mercury's orbit. Since approximate values for the inclination and the node had already been obtained, these elements, which had only a small \SourceLeaf{565}effect on the reduction, could safely be left out at first. The differences between the observed and computed longitudes were expressed as a function of the corrections to the epoch, the mean motion, the eccentricity, the aphelion, and the mass of Venus. This gave 117 equations of condition, from which the corrections to the elements taken as the starting point were derived by the method of least squares. This also yielded the third equation for the value of the mass of Venus. Forming the equations of condition from the geocentric longitudes required care and precision. At the greatest elongations, where Mercury has chiefly and almost exclusively been observed, reduction to the geocentric position greatly diminishes the heliocentric error. Small errors in the observed geocentric longitudes therefore greatly alter the corrections to the elements that are being sought. The author believes that the number of observations eliminated this effect of observational errors and that correcting the solar tables each time from contemporaneous observations avoided the effect of the Sun's positions on the computed longitudes.

With these improved elements the observed geocentric latitudes were reduced to heliocentric latitudes. Comparing them with the computed latitudes gave a hundred equations of condition, which yielded the inclination and the node. The hope of determining the inclination for an earlier epoch and thus obtaining a fourth equation for the mass of Venus from the observed secular change in the inclination was frustrated by the lack of suitable observations.

The final results of these investigations were as follows.
If the mass of Venus as adopted by \textsc{Laplace} (Mécan. Céleste T.\,3. p.\,61) is set $= 1$, the improved mass, as it follows from the three equations mentioned, is $= 1{,}0974 = \frac{1}{849132}$ of the Sun's mass. All periodic and secular perturbations were calculated with this mass.

\begin{center}\fontsize{9}{11}\selectfont\setlength{\tabcolsep}{2pt}\renewcommand{\arraystretch}{1.18}\begin{tabular}{@{}l c r@{\;}r@{\;}r@{}}
Epoch 1750, Seeberg meridian & $.\;.\;.\;.\;.\;.$ & $253^\circ$ & $5'$ & $17''1$\\
Mean annual motion & $.\;.\;.\;.\;.\;.\;.\;.$ & $53$ & $43$ & $3.613$\\
Aphelion 1750 & $.\;.\;.\;.\;.\;.\;.\;.\;.\;.\;.\;.\;.$ & $253$ & $33$ & $24.3$\\
Eccentricity 1800 & $.\;.\;.\;.\;.\;.\;.\;.\;.\;.\;.\;.$ & \multicolumn{3}{l}{$0.2056163$}\\
Semimajor axis & $.\;.\;.\;.\;.\;.\;.\;.\;.\;.\;.\;.$ & \multicolumn{3}{l}{$0.3870988$}\\
Node 1750 & $.\;.\;.\;.\;.\;.\;.\;.\;.\;.\;.\;.\;.\;.\;.$ & $75^\circ$ & $22'$ & $0''96$\\
Inclination of the orbit 1800 & $.\;.\;.\;.\;.\;.\;.\;.\;.\;.$ & $7$ & $0$ & $5.9$\\
Secular change in the aphelion & $.\;.\;.\;.\;.\;.\;.$ & $+1$ & $33$ & $22.9$\\
\quad —\quad —\quad in the node & $.\;.\;.\;.\;.\;.\;.\;.$ & $+1$ & $10$ & $15.1$\\
\quad —\quad —\quad in the eccentricity & $.\;.\;.\;.\;.\;.$ & \multicolumn{3}{l}{$+0''791$}\\
\quad —\quad —\quad in the inclination & $.\;.\;.\;.\;.\;.\;.\;.$ & \multicolumn{3}{l}{$+18.380$}\\
\end{tabular}\end{center}

{\rightskip=0pt plus 3em These elements represent all one hundred observed geocentric positions included in the equations of condition extremely satisfactorily. The same is true of fifty more recent observations of Mercury made by \textsc{Bouvard} in Paris and by the author himself at the Seeberg observatory, which were not used to establish the elements. There is thus every reason to hope that these determinations will continue to represent Mercury's motions in agreement with the sky in the coming decades.\par}
